
%BeginMC
\question

\begin{choices}
\correctchoice  %EndChoice
\choice  %EndChoice
\choice  %EndChoice
\choice  %EndChoice
\choice  %EndChoice
\end{choices}
%EndMC

%%%%%%%%%%%% Questions begin here. Copy and paste from template above.

%%% Sea otter case study

%BeginMC
\question
After transient orca began to prey on sea otters, what happened to the kelp forest?

\begin{choices}
	\choice The kelp abundance remained about the same. The kelp are not part of this community, so the kelp could not be affected.%EndChoice
	\choice The kelp abundance remained about the same. Only the predator-prey interactions between Orca and otters were important, so nothing happened to the kelp.%EndChoice
	\choice The kelp abundance increased. Otters dislodge kelp with their playful antics.  Because there were fewer otters, less kelp was dislodged so the kelp abundance increased.%EndChoice
	\correctchoice The kelp abundance decreased. Sea otters eat urchins and urchins eat kelp.  Because there were fewer otters to eat urchins, the number of urchins increased.  The urchins therefore ate more kelp so the kelp abundance decreased.%EndChoice
	\choice The kelp abundance decreased. Sea otters and kelp have a symbiotic relationship. Without the sea otters, the symbiotic relationship was broken and so the kelp abundance decreased.%EndChoice
\end{choices}
%EndMC

%BeginMC
\question
If humans stopped fishing along the coast of Alaska, what will probably happen to the kelp forest and why?

\begin{choices}
	\correctchoice The kelp abundance will increase because the sea urchin abundance will probably decrease.%EndChoice
	\choice The kelp abundance will remain about the same. The kelp are not part of this community, so the kelp would not be affected.%EndChoice
	\choice The kelp abundance will decrease because more fish will be present to eat the kelp.%EndChoice
	\choice The kelp abundance will decrease. Sea otters eat urchins and urchins eat kelp. Because there were fewer otters to eat urchins, the number of urchins increased. The urchins therefore ate more kelp so the kelp abundance decreased.%EndChoice
	\choice The kelp abundance will increase because of the symbiotic relationship between kelp and fishes. Fishes are necessary for kelp to reproduce.%EndChoice
\end{choices}
%EndMC

%BeginMC
\question
Which of these species is most likely a keystone species in Alaska?

\begin{choices}
	\choice Humans.  %EndChoice
	\choice Kelp.  %EndChoice
	\choice Killer whales. %EndChoice
	\correctchoice Sea otters. %EndChoice
	\choice Sea urchins.  %EndChoice
\end{choices}
%EndMC


%%% Ecological Succesion %%%

%BeginMC
\question
Early successional species modify the environment, which allows intermediate successional species to colonize. This is known as

\begin{choices}
\correctchoice facilitation. %EndChoice
\choice tolerance. %EndChoice
\choice inhibition %EndChoice
\choice secondary succession %EndChoice
\choice primary succession %EndChoice
\choice disturbance %EndChoice
\end{choices}
%EndMC

%BeginMC
\question
Intermediate successional species modify the environment, which allows late successional species to colonize. This is known as

\begin{choices}
	\correctchoice facilitation. %EndChoice
	\choice tolerance. %EndChoice
	\choice inhibition %EndChoice
	\choice secondary succession %EndChoice
	\choice primary succession %EndChoice
	\choice tertiary succession %EndChoice
	\choice disturbance %EndChoice
\end{choices}
%EndMC

%%%% Climate, Climographs, and Ecosystems %%%

%BeginMC
\question
Lower elevations on a very high mountain near the equator experience more precipitation and warmer temperatures than high elevations. If you climbed from the base of the mountain to the top, you would notice a change in the types of ecosystems present on the mountain side. This would be similar to the changes (Note: The equator is 0°, the poles are at 90°). 

\begin{choices}
\choice traveling across South American from 55°S latitude to 5°S latitude.%EndChoice
\correctchoice traveling across North America from 15°N latitude to 55°N latitude.%EndChoice
\choice traveling across South America and North American from 55°S latitude to 55°N latitude.%EndChoice
\choice across Europe and Asia from 15°E longitude to 15°W longitude.%EndChoice
\choice across Europe and Asia from 15°W longitude to 15°E longitude.%EndChoice
\end{choices}
%EndMC

%BeginMC
\question
Probably the most important factor(s) affecting the distribution of ecosystems is (are)

\begin{choices}
\correctchoice climate.%EndChoice
\choice wind and ocean water current patterns.%EndChoice
\choice wind direction and rain shadows.%EndChoice
\choice longitude and latitude.%EndChoice
\choice day length and rainfall.%EndChoice
\end{choices}
%EndMC


%BeginMC
\question
The Mediterranean coast and coastal California have similar types of ecosystems. This is because

\begin{choices}
\correctchoice they have similar climates.%EndChoice
\choice they have similar types of organisms.%EndChoice
\choice they are both wine-growing regions.%EndChoice
\choice they have similar types of communities.%EndChoice
\choice they are at similar latitudes.%EndChoice
\end{choices}
%EndMC


%BeginMC
\fullwidth{\centering\includegraphics[width=0.9\textwidth]{climograph}}
\question
Boreal forests in Canada contain large evergreen trees like spruce trees and pine trees. If a boreal forest continues to get about 140 cm of precipitation annually but mean annual temperature increases from 0°C to 5°C over the next 100 years, the ecosystem will

\begin{choices}
\choice become more like a sclerophyllous woodland. %EndChoice
\choice remain a boreal forest. %EndChoice
\choice become more like a tundra. %EndChoice
\choice become more like a temperate rain forest. %EndChoice
\correctchoice become more like a temperate deciduous forest. %EndChoice
\end{choices}
%EndMC

%BeginMC
\fullwidth{\centering\includegraphics[width=0.9\textwidth]{climograph}}
\question
According to this climograph, the temperate grassland 

\begin{choices}
\choice never occurs with deserts. %EndChoice
\choice has an average annual rainfall of about 20--80 cm. %EndChoice
\choice has an average annual rainfall of about 5--30 cm. %EndChoice
\correctchoice has an average annual temperature range of about 4--30°C. %EndChoice
\choice has the greatest variation in precipitation of any ecosystem shown in the figure. %EndChoice
\end{choices}
%EndMC


%BeginMC
\question
The southeastern region of Missouri lies within the temperate broadleaf forest ecosystem. According to the climograph below, the temperate broadleaf forest

\begin{multicols}{2} 
\begin{choices}
\choice never has more than about 220 cm (87 inches) of rainfall per year.%EndChoice
\correctchoice has an average of 60--220 cm of precipitation per year.%EndChoice
\choice has more rain and similar temperatures to the northern coniferous forest.%EndChoice
\choice always has more precipitation than the temperate grassland.%EndChoice
\choice is never warmer than the temperate grassland.%EndChoice
\end{choices}
\columnbreak
	\hfill \includegraphics[width=6cm]{img/climograph.png}
\end{multicols}
%EndMC


%BeginMC
\fullwidth{\centering\includegraphics[width=0.9\textwidth]{climograph}}
\question
The Ozark region of Missouri lies within the temperate deciduous forest ecosystem.  If, within the next 25 years, global climate change raises the mean annual temperature from 14°C to 18°C and lowers the mean annual precipitation from 160 cm to 130 cm, which of the following is most likely?

\begin{choices}
\choice The ecosystem will remain temperature deciduous ecosystem. %EndChoice
\choice The ecosystem will become more like boreal forest. %EndChoice
\choice The ecosystem will become more like temperate rain forest. %EndChoice
\choice The ecosystem will become more like temperate grassland. %EndChoice
\correctchoice The ecosystem will become more like sclerophyllous or thorn woodland. %EndChoice
\end{choices}
%EndMC


%BeginMC
\question
The Ozark region of Missouri lies within the temperate deciduous forest ecosystem.  If, within the next 25 years, global climate change raises the mean annual temperature from 14°C to 18°C and lowers the mean annual precipitation from 160 cm to 130 cm, which of the following is most likely?

\begin{multicols}{2}
	\begin{choices}
		\choice The ecosystem will remain temperature deciduous ecosystem. %EndChoice
		\choice The ecosystem will become more like boreal forest. %EndChoice
		\choice The ecosystem will become more like temperate rain forest. %EndChoice
		\choice The ecosystem will become more like temperate grassland. %EndChoice
		\correctchoice The ecosystem will become more like sclerophyllous or thorn woodland. %EndChoice
	\end{choices}
\columnbreak
	\includegraphics[width=0.45\textwidth]{exam4_climograph2}
\end{multicols}
%EndMC

%BeginMC
\question
The African savanna is famous for its wildlife, such as lions and giraffes. Based on the climograph, which of the following most accurately defines the savanna ecosystem?

\begin{choices}
\choice Savannas receive lots of sunlight and moderate amounts of rain. %EndChoice
\correctchoice Savannas have mean annual precipitation of about 75--130 cm and mean annual temperatures of about 18--29°C. %EndChoice
\choice Savannas are much drier than temperate grasslands but have about the same annual mean temperature.  %EndChoice
\choice Savannas are drier than thorn scrub ecosystems but have about the same temperatures. %EndChoice
\choice A typical savanna might have about a mean annual rainfall of about 25 cm and mean annual temperature of about 100°C. %EndChoice
\end{choices}
%EndMC


%BeginMC
\fullwidth{\centering\includegraphics[width=0.9\textwidth]{climograph}}
\question
A town gets an average annual precipitation of 102 cm per year. The average annual temperature is 14°C. This town is probably located in 

\begin{choices}
\correctchoice a temperate grassland ecosystem. %EndChoice
\choice a savannah ecosystem. %EndChoice
\choice a sclerophyllous woodland. %EndChoice
\choice a scrubland ecosystem. %EndChoice
\choice a temperate deciduous forest ecosystem. %EndChoice
\choice a thorn woodland ecosystem. %EndChoice
\end{choices}
%EndMC



%BeginMC
\question
Higher elevations on a mountain experience less precipitation and cooler temperatures than low elevations.  If you hiked dowm from the top of the mountain, you would notice a change in the types of plants and animals present in the community.  This is analogous to the changes in the ecosystems if you walked 

\begin{choices}
\correctchoice from the north pole to the equator.%EndChoice
\choice east to west over different altitudes.%EndChoice
\choice west to east at one latitude.%EndChoice
\choice from the equator to the south pole.%EndChoice
\choice from the north pole to the south pole.%EndChoice
\end{choices}
%EndMC


%%% Rain Shadow %%%

%BeginMC
\question
The downwind side of large mountain ranges is often drier than the upwind side. This phenomenon occurs because

\begin{choices}
\choice mountains are always near oceans.%EndChoice
\choice because the rainforest on the west side of the mountain absorbs all of the rain out of the air before it crosses the mountain.%EndChoice
\choice when the desert on the east side of the mountains sucks the moisture out of the air.%EndChoice
\correctchoice the mountain range deflects the moisture-laden air from the ocean up into the atmosphere. The air cools and drops its moisture as rain. The now-dry air flows across downwind side of the mountains.%EndChoice
\choice the wind blows the moisture out of the air before it can cross the mountains.%EndChoice
\end{choices}
%EndMC



%BeginMC
\question
The Atacama Desert, said to be the driest desert on earth, gets an average of 15 mm (0.6 in) of precipitation a year.  The desert is about 8000 m (16 miles!) above sea level. The Atacama Desert is probably drier than other deserts because

\begin{choices}
\choice the Atacama is warmer than most other deserts because it is closer to the sun (because it is so high). %EndChoice
\choice the Atacama is colder than most other deserts because it is so high. %EndChoice
\correctchoice the Atacama is made even drier by the rain shadow of the Andes Mountains. %EndChoice
\choice the Atacama receives more sunlight than most other deserts, causing more evaporation. %EndChoice
\choice the Atacama has an average annual temperature above 2.5°C and an average annual precipitation below 50 cm. %EndChoice
\end{choices}
%EndMC

%BeginMC
\question
A rain shadow results in 

\begin{choices}
\correctchoice desert-like conditions on the downwind side of a mountain due. %EndChoice
\choice extremely wet conditions due to the lack of wind over a mountain range. %EndChoice
\choice dry air moving towards the north pole that cools and sinks, creating Arctic tundra. %EndChoice
\choice dry air near the top of the mountains, creating alpine tundra. %EndChoice
\choice global polar regions that rarely receive moisture from the warmer, tropical regions, and are therefore dryer. %EndChoice
\end{choices}
%EndMC

%BeginMC
\question
The land on the south side of the Himalayan mountain range (below) receives relatively high levels of rainfall and has fertile farmlands.  The land north of the mountain range is hot and dry and can't be farmed.  Which of the following is true.

\begin{multicols}{2}
\begin{choices}
\choice The south side of the Himalayan mountain range is in a rain shadow.%EndChoice
\choice The desert in southern China warms the north side of the mountain range, causing it to dry out.%EndChoice
\choice The cool, moist Indian Ocean air heats up as it rises, releasing its precipitation as it passes the tops of the mountains, and this warm, now dry air cools as it descends on the leeward side of the range.%EndChoice
\correctchoice Warm, moist air from the Indian Ocean blows north, rises and cools, releasing precipitation as it moves up the south side of the range, and this cool, now dry air mass heats up as it descends on the north side of the range.%EndChoice
\choice The warm, moist air is too dense to rise over the height of Mount Everest and the other mountains in this range so the moisture remains on the south side of the range.%EndChoice
\end{choices}
\columnbreak
	\includegraphics[width=8cm]{img/himalayas.png}
\end{multicols}
%EndMC


%BeginMC
\question
Air masses formed over the Pacific Ocean are moved by prevailing westerlies where they encounter extensive north-south mountain ranges, such as the Sierra Nevada and the Cascades. Which statement best describes the changes that these air masses undergo?  

\begin{choices}
\correctchoice The warm, moist Pacific air rises and cools, releasing precipitation as it moves up the windward side of the range, and this cool, now dry air mass heats up as it descends on the leeward side of the range.  %EndChoice
\choice The cool, moist Pacific air heats up as it rises, releasing its precipitation as it passes the tops of the mountains, and this warm, now dry air cools as it descends on the leeward side of the range. %EndChoice
\choice The cool, dry Pacific air heats and picks up moisture from evaporation of the snowcapped peaks of the mountain range, releasing this moisture as precipitation as the air cools as it descends on the leeward side of the range.  %EndChoice
\choice These air masses are blocked by these mountain ranges producing high annual amounts of precipitation on the windward sides of these mountain ranges.  %EndChoice
\choice These air masses remain essentially unchanged in moisture content and temperature as they pass over these mountain ranges.  %EndChoice
\end{choices}
%EndMC


%%% Trophic Levels and Cascades %%



%BeginMC
\question
In general, more herbivores will be found in an ecosystem than carnivores. Why?

\begin{choices}
\correctchoice There is more total energy at the primary consumer trophic level to support more biomass.%EndChoice
\choice There is more total energy at the secondary consumer trophic level to support larger but fewer carnivores.%EndChoice
\choice More energy is transferred from primary producers to primary consumers than from primary consumers to secondary consumers.%EndChoice
\choice Herbivores are usually faster than carnivores so they survive longer and reproduce more often.%EndChoice
\choice More energy is transferred from plants to herbivores than from herbivores to carnivores.%EndChoice
\end{choices}
%EndMC


%BeginMC
\question
On average, the amount of primary consumer biomass is about \underline{\hspace{2cm}} less than the amount of primary producer biomass.

\begin{choices}
\choice 10\% %dontshuffle %EndChoice
\choice 20\% %dontshuffle %EndChoice
\choice 50\% %dontshuffle %EndChoice
\choice 80\% %dontshuffle %EndChoice
\correctchoice 90\% %dontshuffle %EndChoice
\end{choices}
%EndMC


%BeginMC
\question
On average, the amount of secondary consumer biomass is about \underline{\hspace{2cm}} less than the amount of primary producer biomass. 
\begin{choices}
\choice 10$\times$ %dontshuffle %EndChoice
\choice 20$\times$ %dontshuffle %EndChoice
\choice 50$\times$ %dontshuffle %EndChoice
\choice 80$\times$ %dontshuffle %EndChoice
\correctchoice 100$\times$ %dontshuffle %EndChoice
\end{choices}
%EndMC

%BeginMC
\question
On average, the amount of secondary consumer biomass is about \underline{\hspace{2cm}} less than the amount of primary producer biomass. 
\begin{choices}
\choice 10\times %dontshuffle %EndChoice
\choice 20\times %dontshuffle %EndChoice
\choice 50\times %dontshuffle %EndChoice
\choice 80\times %dontshuffle %EndChoice
\correctchoice 100\times %dontshuffle %EndChoice
\end{choices}
%EndMC


%BeginMC
\question
A transient killer whale eats seals. The seals ate fish. The fish ate part of a dead killer whale laying on the sea floor, but also ate smaller live squid. You do not know what the squid might have eaten. The fish in this example is acting as \underline{\hspace{2cm}}. (Choose the best answer)

\begin{choices}
\correctchoice a consumer and a decomposer.%EndChoice
\choice a secondary consumer and a primary consumer.%EndChoice
\choice a secondary consumer and a decomposer.%EndChoice
\choice a secondary consumer, a decomposer, and a consumer one trophic level higher than the killer whale.%EndChoice
\choice a consumer and a producer.%EndChoice
\end{choices}
%EndMC

%BeginMC
\question
Energy from the sun is converted into chemical energy for organisms by

\begin{choices}
\correctchoice producers.%EndChoice
\choice carnivores.%EndChoice
\choice decomposers.%EndChoice
\choice herbivores.%EndChoice
\end{choices}
%EndMC


%BeginMC
\question
Most human are omnivores (eating plants and animals) but eating a vegetarian diet is considered to be more globally sustainable than eating animal products. Why is this?
 
\begin{choices}
\choice Plants have more food calories per unit of mass than do animal products. %EndChoice
\choice Energy is much more concentrated in plant materials than in animals. %EndChoice
\choice Toxins become more concentrated in the higher levels of the ecological pyramid. %EndChoice
\choice More energy is transferred from plants to animals than from animals to animals. %EndChoice
\correctchoice Vegetarianism is more ecologically energy efficient than eating animal products. %EndChoice
\end{choices}
%EndMC

%BeginMC
\question
The energy that is \textit{not} transferred to other trophic levels or decomposers is lost

\begin{choices}
\correctchoice as metabolic heat.%EndChoice
\choice in fecal matter.%EndChoice
\choice to consumers at any trophic level.%EndChoice
\choice to primary consumers.%EndChoice
\choice to secondary consumers.%EndChoice
\end{choices}
%EndMC


%BeginMC
\question
The number of herbivores found in an trophic pyramid is limited by the

\begin{choices}
\choice number of producers below the herbivores.%EndChoice
\choice number of predators that eat the herbivores.%EndChoice
\correctchoice amount of energy transferred from producers to the next trophic level.%EndChoice
\choice amount of energy in the entire pyramid.%EndChoice
\choice number of trophic levels (e.g., 3 or 4 levels) above the herbivores.%EndChoice
\end{choices}
%EndMC


%BeginMC
\question
The number of carnivores found at the top of a trophic pyramid is limited by the

\begin{choices}
\choice number of organisms below the top carnivores. %EndChoice
\choice number of trophic levels below the top carnivores. %EndChoice
\correctchoice amount of energy transferred to the top carnivores. %EndChoice
\choice amount of biomass below the top carnivores. %EndChoice
\choice amount of energy recycled by decomposers on the next lower trophic level. %EndChoice
\end{choices}
%EndMC

%BeginMC
\question
What type of organisms capture sunlight and convert it into forms of energy than can be used by other organisms in an ecosystem?

\begin{choices}
\choice Herbivores %EndChoice
\choice Carnivores %EndChoice
\correctchoice Producers %EndChoice
\choice Detritivores %EndChoice
\choice Decomposers %EndChoice
\end{choices}
%EndMC

%BeginMC
\question
A fish ate shrimp. The shrimp ate algae (which are photosynthetic).  The shrimp is a \rule{1in}{0.4pt}. The algae is a \rule{1in}{0.4pt}.

\begin{choices}	
\choice primary consumer; secondary consumer. %EndChoice
\choice secondary consumer; primary producer. %EndChoice
\choice primary consumer; decomposer. %EndChoice
\choice decomposer; primary producer. %EndChoice
\correctchoice primary consumer; primary producer. %EndChoice
\end{choices}
%EndMC

%BeginMC
\question
A fish ate shrimp. The shrimp ate algae (which are photosynthetic).  The shrimp is a \rule{1in}{0.4pt}. The fish is a \rule{1in}{0.4pt}.

\begin{choices}	
\correctchoice primary consumer; secondary consumer. %EndChoice
\choice secondary consumer; primary producer. %EndChoice
\choice primary consumer; decomposer. %EndChoice
\choice decomposer; primary producer. %EndChoice
\choice primary consumer; primary producer. %EndChoice
\end{choices}
%EndMC

%BeginMC
\question
A fish ate shrimp. The shrimp ate algae (which are photosynthetic).  The algae is a \rule{1in}{0.4pt}. The fish is a \rule{1in}{0.4pt}.

\begin{choices}	
\choice primary consumer; secondary consumer. %EndChoice
\choice secondary consumer; primary consumer. %EndChoice
\choice secondary consumer; primary producer. %EndChoice
\choice primary consumer; primary producer. %EndChoice
\correctchoice primary producer; secondary consumer. %EndChoice
\choice primary producer; primary consumer. %EndChoice
\end{choices}
%EndMC

%BeginMC
\question
A transient killer whale eats sea otters. The sea otters ate sea urchins. The sea urchins ate live kelp from the sea floor.  The sea otter is a \rule{1in}{0.4pt}. The sea urchin is a \rule{1in}{0.4pt}.	

\begin{choices}
\choice primary consumer; secondary consumer. %EndChoice
\correctchoice secondary consumer; primary consumer. %EndChoice
\choice primary consumer; primary producer. %EndChoice
\choice secondary consumer; decomposer. %EndChoice
\choice tertiary consumer; secondary consumer. %EndChoice
\end{choices}
%EndMC


%BeginMC
\question
Squid eat fish.  Fish eat shrimp. Shrimp eat algae.  If algae has 10,000 Joules of energy, then how much energy makes it to squid.

\begin{choices}	
\correctchoice 10 Joules. %EndChoice
\choice 60 Joules. %EndChoice
\choice 100 Joules. %EndChoice
\choice 1000 Joules. %EndChoice
\choice 600 Joules. %EndChoice
\end{choices}
%EndMC


%BeginMC
\question
Organisms that break down organic material and make it available to other organisms are called

\begin{choices}
\correctchoice decomposers. %EndChoice
\choice scavengers. %EndChoice
\choice consumers.  %EndChoice
\choice producers. %EndChoice
\choice degraders. %EndChoice
\end{choices}
%EndMC


%BeginMC
\question
A tertiary consumer has high abundance. In a trophic cascade, which of the following is probably \emph{not} true.
\begin{choices}
	\correctchoice The primary producer will have high abundance. %EndChoice
	\choice The secondary consumer will have low abundance. %EndChoice
	\choice The primary consumer will have high abundance. %EndChoice
	\choice The decomposer level will not be affected. %EndChoice%
	\choice The quaternary consumer (next level higher) will have high abundance. %EndChoice
\end{choices}
%EndMC

%BeginMC
\question
A secondary consumer has high abundance. In a trophic cascade, which of the following is probably \emph{not} true.
\begin{choices}
\correctchoice The primary producer will have low abundance. %EndChoice
\choice The tertiary consumer will have low abundance. %EndChoice
\choice The primary consumer will have low abundance. %EndChoice
\choice The different trophic levels will have alternating high and low abundances. %EndChoice
\choice Actually, all of the above are true. %dontshuffle %EndChoice%
\end{choices}
%EndMC

%BeginMC
\question
A field containing 2000 kg of plant material can support about \rule{1in}{0.4pt} of tertiary consumers.

\begin{choices}
\correctchoice 2 kg %EndChoice
\choice 20 kg %EndChoice
\choice 10 kg %EndChoice
\choice 10\% %EndChoice
\choice 200 kg %EndChoice
\choice 90\% %EndChoice
\end{choices}
%EndMC

%BeginMC
\question
Based on the hypothetical food web below, which letter most likely represents a decomposer. Arrows point in the direction of energy flow.

\begin{multicols}{3} 
	\begin{choices}
		\correctchoice A %dontshuffle %EndChoice
		\choice B %dontshuffle %EndChoice
		\choice C %dontshuffle %EndChoice
		\choice D %dontshuffle %EndChoice
		\correctchoice E %dontshuffle %EndChoice
	\end{choices}
	\columnbreak
	\includegraphics{exam4_foodweb}
	\columnbreak
\end{multicols}
%EndMC

%BeginMC
\question
Based on the hypothetical food web below, one letter represents an organism that could be both a secondary and tertiary consumer? The arrows point in the direction of energy transfer from lower to higher trophic levels. More than one answer is possible. Circle only one answer.

\begin{multicols}{3} 
	\begin{choices}
		\correctchoice A %dontshuffle %EndChoice
		\choice B %dontshuffle %EndChoice
		\choice C %dontshuffle %EndChoice
		\choice D %dontshuffle %EndChoice
		\correctchoice E %dontshuffle %EndChoice
	\end{choices}
\columnbreak
	\includegraphics{exam4_foodweb}
\columnbreak
\end{multicols}
\vspace{-1\baselineskip}
%EndMC

%% Same as above question but using trophic.
%BeginMC
\question
Based on the hypothetical food web below, one letter represents an organism that could be both a secondary and tertiary consumer? The arrows point from lower to higher trophic levels. More than one answer is possible. Circle only one answer.

\begin{multicols}{3} 
	\begin{choices}
		\correctchoice A %dontshuffle %EndChoice
		\choice B %dontshuffle %EndChoice
		\choice C %dontshuffle %EndChoice
		\choice D %dontshuffle %EndChoice
		\correctchoice E %dontshuffle %EndChoice
	\end{choices}
\columnbreak
	\includegraphics{exam4_foodweb}
\columnbreak
\end{multicols}
\vspace{-1\baselineskip}
%EndMC
 
%BeginMC
\question
Based on the hypothetical food web below, which letter represents an organism that is most likely a decomposer?  The arrows point to organisms receiving energy.

\begin{multicols}{3} 
	\begin{choices}
		\correctchoice A %dontshuffle %EndChoice
		\choice B %dontshuffle %EndChoice
		\choice C %dontshuffle %EndChoice
		\choice D %dontshuffle %EndChoice
		\choice E %dontshuffle %EndChoice
	\end{choices}
\columnbreak
	\includegraphics{exam4_foodweb}
\columnbreak
\end{multicols}
\vspace{-1\baselineskip}
%EndMC
 
 %BeginMC
\question
Based on the hypothetical food web below, which letter represents an organism that is most likely a primary producer?  The arrows point to organisms receiving energy.

\begin{multicols}{3} 
	\begin{choices}
		\choice A %dontshuffle %EndChoice
		\correctchoice B %dontshuffle %EndChoice
		\choice C %dontshuffle %EndChoice
		\choice D %dontshuffle %EndChoice
		\choice E %dontshuffle %EndChoice
	\end{choices}
\columnbreak
	\includegraphics{exam4_foodweb}
\columnbreak
\end{multicols}
\vspace{-1\baselineskip}
%EndMC
 
%BeginMC
\question
Consider the food chain grass $\to$ grasshopper $\to$ mouse $\to$ snake $\to$ hawk. How much of the chemical energy fixed by photosynthesis of the grass (100\%) is available to the hawk?

\begin{choices}
\correctchoice 0.01\% %dontshuffle %EndChoice
\choice 0.1\% %dontshuffle %EndChoice
\choice 1\% %dontshuffle %EndChoice
\choice 10\% %dontshuffle %EndChoice
\choice 60\% %dontshuffle %EndChoice
\end{choices}
%EndMC





%%% General Ecosystems %%%


%BeginMC
\question
Ecologists say that energy ``flows'' through an ecosystem because

\begin{choices}
\choice energy, like air, can be thought of as a fluid that moves through the ecosystem.%EndChoice
\correctchoice the heat energy that is created by metabolism can never again enter the ecosystem.%EndChoice
\choice energy flows from high trophic levels, such as top predators, to low trophic levels, such as prey.%EndChoice
\choice energy cycles from one trophic level to the next.%EndChoice
\choice all energy that enters an ecosystem flows there from the sun.%EndChoice
\end{choices}
%EndMC


%BeginMC
\question
The major abiotic (non-living) components of ecosystems include

\begin{choices}
\correctchoice temperature, moisture, sunlight, and nutrients.%EndChoice
\choice plants, animals, climate, and microorganisms.%EndChoice
\choice producers, herbivores, carnivores, and plants.%EndChoice
\choice consumers, detritus feeders, decomposers, and parasites.%EndChoice
\choice producers, consumers, decomposers, and detritus feeders.%EndChoice
\end{choices}
%EndMC



%BeginMC
\question
Plants produce water vapor during photosynthesis, which goes from the plant to the atmosphere. The water vapor in the atmosphere eventually falls back to the earth as rain. This process would be part of

\begin{choices}
\correctchoice the water cycle.%EndChoice
\choice energy flow.%EndChoice
\choice plant metabolism.%EndChoice
\choice the entire ecosystem.%EndChoice
\choice the energy cycle.%EndChoice
\choice the nutrient cycle.%EndChoice
\end{choices}
%EndMC

%BeginMC
\question
Which of the following is an example of an ecosystem?

\begin{choices}
\choice All of the brook trout in a 500 hectare\textsuperscript{2} river drainage system.%EndChoice
\correctchoice The plants, animals and decomposers, plus the nutrients and rainfall in a tropical rain forest.%EndChoice
\choice The plants, animals, and decomposers that inhabit an alpine meadow.%EndChoice
\choice A pond and all of the plant and animal species that live in it.%EndChoice
\choice The intricate interactions of the various plant and animal species on a savanna during a drought.%EndChoice
\end{choices}
%EndMC


%BeginMC
\question
Which of the following is a complete example of an ecosystem?  

\begin{choices}
\choice All of the animals on a coral reef. %EndChoice
\choice All of the plants, animals, decomposers and detritivores in a prairie grassland. %EndChoice
\choice A pond and all the animal species that live in it, plus the sunlight that hits the surface of the pond. %EndChoice
\choice The transfer of energy from primary producer to the top level consumer. %EndChoice
\correctchoice Interactions between producers, consumers, and decomposers with each other and with sunlight, nutrients, water and other chemical elements. %EndChoice
\end{choices}
%EndMC


%BeginMC
\question
Which statement most accurately describes how nutrients and energy are used in ecosystems?

\begin{choices}
\choice Nutrients flow through ecosystems; energy cycles within ecosystems. %EndChoice
\correctchoice Energy flows through ecosystems; nutrients cycle within ecosystems. %EndChoice
\choice Energy can be converted into nutrients; nutrients cannot be converted into energy. %EndChoice
\choice Nutrients can be converted into energy; energy cannot be converted into nutrients. %EndChoice
\choice Nutrients are used in ecosystems; energy is not. %EndChoice
\end{choices}
%EndMC


%BeginMC
\question
Which of the following statements about ecosystems is false?

\begin{choices}
\correctchoice On average, 10\% of the energy available at one trophic is lost as it moves to the next higher trophic level. %EndChoice
\choice Energy flows through an ecosystem. %EndChoice
\choice Nutrients, like nitrogen and phosphorus, cycle within an ecosystem. %EndChoice
\choice Water cycles within an ecosystem. %EndChoice
\choice Carnivores occupy higher trophic levels than herbivores. %EndChoice
\end{choices}
%EndMC


%BeginMC
\question
The energy stored by producers can be passed through an ecosystem along a(n) \rule{1in}{0.4pt}, a series of interactions in which organisms transfer energy by eating and being eaten.

\begin{choices}
\correctchoice food web %EndChoice
\choice ecosystem %EndChoice
\choice trophic cascade %EndChoice
\choice energy cycle %EndChoice
\choice nutrient cycle %EndChoice
\choice predator-prey chain  %EndChoice
\end{choices}
%EndMC



%BeginMC
\question
All energy that enters an ecosystem from the sun will eventually end up as

\begin{choices}
\correctchoice heat energy. %EndChoice
\choice fecal material. %EndChoice
\choice nutrients. %EndChoice
\choice chemical energy. %EndChoice
\choice cellular respiration. %EndChoice
\choice growth for organisms. %EndChoice
\end{choices}
%EndMC

%BeginMC
\question
Probably the most important factor affecting the global distribution of ecosystems are

\begin{choices}
\choice average annual wind patterns and daily ocean current patterns. %EndChoice
\choice daily wind patterns and average annual precipitation. %EndChoice
\choice day length and average annual rainfall. %EndChoice
\choice average annual rainfall and average annual temperature. %EndChoice
\correctchoice average annual precipitation and average annual temperature. %EndChoice
\end{choices}
%EndMC


%BeginMC
\question
Which of the following areas of study focuses on the exchange of energy and nutrients among individuals of different species?

\begin{choices}
\correctchoice Ecosystem ecology. %EndChoice
\choice Population Ecology %EndChoice
\choice Evolutionary Ecology %EndChoice
\choice Community Ecology %EndChoice
\choice Behavioral Ecology %EndChoice
\end{choices}
%EndMC

%BeginMC
\question
In ecosystems, why is ``cycling'' used to describe the transfer of chemical elements and nutrients, while ``flow'' is used for the transfer of energy?

\begin{choices}
\correctchoice Chemical elements are recycled within ecosystems, but energy flows through and out of ecosystems.%EndChoice
\choice Both chemical elements and energy are recycled but can flow to other ecosystems.%EndChoice
\choice Chemical elements are cycled into ecosystems from other ecosystems, but energy constantly flows within the ecosystem.%EndChoice
\choice Both chemical elements and energy flow within an ecosystem but can cycle to other ecosystems.%EndChoice
\choice Energy flows from one ecosystem to other ecosystems, but chemical elements constantly cycle within an ecosystem.%EndChoice
\end{choices}
%EndMC

%BeginMC
\question
Which of the following are important biotic (living) factors that can affect the  organization of communities?

\begin{choices}
\choice precipitation, wind%EndChoice
\choice nutrient availability, soil pH%EndChoice
\correctchoice predation, competition%EndChoice
\choice temperature, water%EndChoice
\choice light intensity, seasonality%EndChoice
\end{choices}
%EndMC


%%%% Dead Zone %%%

%BeginMC
\question
Considering your knowledge of the dead zone and the causes of microscopic algae blooms, the Gulf of Mexico ecosystem is \emph{most likely} to be limited by which of the following nutrients?

\begin{choices}
	\choice carbon%EndChoice
	\choice \ce{CO2}%EndChoice
	\choice hydrogen%EndChoice
	\choice oxygen%EndChoice
	\correctchoice nitrogen%EndChoice
\end{choices}
%EndMC

%BeginMC
\question
In summer 2012, the Midwest experienced one of the worst droughts in recorded history. Which of the following statements is the best explanation of how this drought affected the Gulf of Mexico dead zone during summer 2012?

\begin{choices}
\correctchoice The dead zone will be smaller because the Mississippi River will carry less water so fewer nutrients will be carried to the Gulf of Mexico.%EndChoice
\choice The dead zone will be smaller because farmers, anticipating the drought, would not have added fertilizer to the crop fields the spring before the drought.%EndChoice
\choice The dead zone will be similar to previous years because the drought is not likely to affect the amount of water that flows out through the Mississippi River.%EndChoice
\choice The dead zone will be larger because the reduced volume of water will concentrate the nutrients that enter the Gulf of Mexico.%EndChoice
\choice The dead zone will be smaller because less freshwater will enter the Gulf of Mexico, allowing more saltwater fishes to live around the mouth of the river.%EndChoice
\end{choices}
%EndMC


%BeginMC
\question
The microscopic algae community in the Gulf of Mexico is limited by which of the following nutrients?

\begin{choices}
\choice carbon%EndChoice
\choice \ce{CO2}%EndChoice
\choice hydrogen%EndChoice
\choice oxygen%EndChoice
\correctchoice nitrogen%EndChoice
\end{choices}
%EndMC


%BeginMC
\question
In 2011, the Mississippi River experienced a record flood event.  Which of the following statements is the best prediction of how this record flood event affected the Gulf of Mexico dead zone that summer?

\begin{choices}
\correctchoice The dead zone will be larger because the volume of water will carry even more nutrients into the Gulf of Mexico.%EndChoice
\choice The dead zone will be smaller because the volume of water will bring more oxygen to the Gulf of Mexico, which will offset the hypoxia.%EndChoice
\choice The dead zone will be smaller because the volume of water will dilute the nutrients, reducing eutrophication.%EndChoice
\choice The dead zone will be similar to previous years because most of the flood water will soak into the ground and not enter the Gulf of Mexico.%EndChoice
\choice The dead zone will be larger because the volume of water will cause more fishes to avoid the freshwater input.%EndChoice
\end{choices}
%EndMC

%BeginMC
\question
In the Gulf of Mexico, eutrophication causes

\begin{choices}
\choice the size of the Dead Zone to increase because of the loss of nutrients.%EndChoice
\correctchoice the community of marine algae to increase due to added nutrients.%EndChoice
\choice the community of marine algae to die due to the loss of nutrients.%EndChoice
\choice the number of shrimp to decrease due to the loss of nutrients.%EndChoice
\choice the number of shrimp to increase due to added nutrients.%EndChoice
\end{choices}
%EndMC

%BeginMC
\question
The dead zone in the Gulf of Mexico and other oceans around the globe occurs when

\begin{choices}
\choice an excess of nutrients causes an algal bloom. The algae consume all of the oxygen in the water for photosynthesis, so other organisms can no longer live there.%EndChoice
\choice eutrophication causes an absence of nutrients in the zone.%EndChoice
\choice commercial fishers over-harvest shrimp, crab, fishes and other sea animals.%EndChoice
\correctchoice an excess of nutrients causes an algal bloom, followed by death of the algae. Decomposition of the algae uses up the oxygen in the water, so other organisms can no longer live there.%EndChoice
\choice global climate change causes the water to become too warm to sustain life.%EndChoice
\end{choices}
%EndMC



%%% Survivorship Curves %%%%


%BeginMC
\question
According to the figure below, which is most correct about the offspring produced by oysters?

\begin{multicols}{2}
\begin{choices}
\correctchoice Most offspring will die at a very young age.%EndChoice
\choice Most offspring will live most of their natural lifespan.%EndChoice
\choice About the same number of offspring will die every year.%EndChoice
\choice Most offspring will die late in life.%EndChoice
\choice Most offspring will die about half way through their natural life.%EndChoice
\end{choices}
\columnbreak
	\hfill \includegraphics[width=6cm]{img/survivorship.png}
\end{multicols}
%EndMC

%BeginMC
\question
According to the figure below, which is most correct about the offspring produced by hydra?

\begin{multicols}{2}
\begin{choices}
\choice Most offspring will die at a very young age.%EndChoice
\choice Most offspring will live most of their natural lifespan.%EndChoice
\correctchoice About the same number of offspring will die every year.%EndChoice
\choice Most offspring will die late in life.%EndChoice
\choice Most offspring will die about half way through their natural life.%EndChoice
\end{choices}
\columnbreak
	\hfill \includegraphics[width=6cm]{img/survivorship.png}
\end{multicols}
%EndMC

%BeginMC
\question
According to the figure below, which is most correct about the offspring produced by Type I organisms?

\begin{multicols}{2}
\begin{choices}
\choice Most offspring will die at a very young age.%EndChoice
\correctchoice Most offspring will live most of their natural lifespan.%EndChoice
\choice About the same number of offspring will die every year.%EndChoice
\choice About 50\% of the individuals will die every year.%EndChoice
\choice Most offspring will die about half way through their natural life.%EndChoice
\end{choices}
\columnbreak
	\hfill \includegraphics[width=6cm]{img/survivorship.png}
\end{multicols}
%EndMC





	